\documentclass[10pt]{article}
\usepackage[margin=0.95in]{geometry}
\usepackage{amsmath,amssymb,amsthm,hyperref}
\newtheorem{theorem}{Theorem}
\newtheorem{lemma}{Lemma}
\newcommand{\T}{\mathbb T}
\newcommand{\G}{\mathcal G}
\title{A shared-reset coupling outside every finite product construction}
\author{Principia Math in Progress}
\date{12 September 2026}
\begin{document}
\maketitle

\paragraph{Task 1.}
MathDB 339681 records Conjecture 19 of N\"usken and Pavliotis
\cite[pp.~337--338]{NP}. For generators of ergodic Feller semigroups,
the conjecture asserts that one can choose decompositions
\[
 L_i=-\sum_{k=1}^{K_i}(A_k^i)^*A_k^i+B^i
\]
so that every coupling operator is a sum
\begin{equation}\label{construction}
 \Gamma=\sum_{\substack{i\ne j\\1\leq k\leq K_i\\1\leq l\leq K_j}}
 \alpha_{ijkl} A_k^i A_l^j ,
\end{equation}
with scalar coefficient functions on the product state space.
We disprove the stated finite-sum construction. The obstruction holds
for every choice of its finite families of scalar operators.

\begin{theorem}
Let each of two marginal processes on the circle \(\T=\mathbb R/\mathbb Z\)
reset, at rate one, to an independent uniform point of the circle.
The coupling that resets both coordinates simultaneously to the same
uniform point is a rich ergodic Feller coupling. Its coupling operator
cannot have the form \eqref{construction} for any finite \(K_1,K_2\)
and any decompositions of the marginal generators.
\end{theorem}

\begin{proof}
Let \(\pi\) be normalized Lebesgue measure on \(\T\), and define
\[
 (Pf)(x)=\int_\T f(z)\,d\pi(z),\qquad L=P-I.
\]
Here \(P\) is a positive, constant-preserving projection of norm one on
\(C(\T)\). Thus
\begin{equation}\label{semigroup}
 S_t=e^{-t}I+(1-e^{-t})P
\end{equation}
is a conservative, uniformly continuous Feller semigroup, with bounded
generator \(L\). It is precisely the stated reset process.
Its invariant measure is uniquely \(\pi\), and \(S_tf\) converges uniformly
to \(\pi(f)\). Ergodic time averages follow directly from the strong law:
successive reset positions are independent with law \(\pi\), and the
independent holding times are exponential with mean one, so the ratio of
the accumulated rewards and holding times converges to \(\pi(f)\).
This applies to every bounded continuous observable and every initial
point. Smooth functions are invariant under \(S_t\) and dense in \(C(\T)\).
Since \(L\) is bounded, they are a core, and the process is rich.

The required marginal decomposition exists. Indeed, \(P\) is an
orthogonal projection in \(L^2(\pi)\), so with \(A=I-P\) and \(B=0\),
\[
 -A^*A=-(I-P)^2=P-I=L.
\]
Our argument below excludes all other finite decompositions as well.

On \(C(\T^2)\), put
\begin{align*}
 (P_xF)(x,y)&=\int_\T F(z,y)\,d\pi(z),&
 (P_yF)(x,y)&=\int_\T F(x,z)\,d\pi(z),\\
 (QF)(x,y)&=\int_\T F(z,z)\,d\pi(z).
\end{align*}
The shared-reset generator and its coupling operator are
\begin{equation}\label{Gamma}
 \overline L=Q-I,\qquad
 \Gamma=\overline L-(L_x+L_y)=Q-P_x-P_y+I.
\end{equation}
The projection \(Q\) is positive and constant-preserving, so
\(\overline S_t=e^{-t}I+(1-e^{-t})Q\) is again a conservative,
uniformly continuous Feller semigroup. Its process resets the pair to
\((Z,Z)\), with \(Z\sim\pi\), at each rate-one Poisson time.
Each coordinate therefore has exactly the prescribed marginal process,
including its path law. The unique invariant distribution is the
diagonal measure \(\Delta_\#\pi\), where \(\Delta(z)=(z,z)\);
the same renewal argument proves ergodicity.

All the operators in \eqref{Gamma} are bounded in the supremum norm,
and \(\Gamma\) is defined on all bounded Borel functions.
It vanishes on functions of either coordinate alone.
On \(C(\T^2)\), \(L_x+L_y+\Gamma=Q-I\) is the bounded Feller generator
just constructed. Smooth functions on \(\T^2\) are a core.
Thus both conditions of the source's Definition 6 hold, including
closability and the generator property. In fact \(\Gamma\in\G^0\).

For any smooth \(f,g\), \eqref{Gamma} gives
\begin{equation}\label{bilinear}
 \Gamma(f\otimes g)(x,y)
 =\int_\T f(z)g(z)\,d\pi(z)
   -g(y)\pi(f)-f(x)\pi(g)+f(x)g(y).
\end{equation}
Now suppose \eqref{construction} holds for finite \(K_1,K_2\).
Operators acting in different coordinates commute on simple tensors.
Combining the terms with \((i,j)=(1,2)\) and \((2,1)\), we obtain
\begin{equation}\label{factor}
 \Gamma(f\otimes g)(x,y)
 =\sum_{k=1}^{K_1}\sum_{l=1}^{K_2}
   \beta_{kl}(x,y)\,(A_k^1f)(x)\,(A_l^2g)(y).
\end{equation}
Consequently, for any \(m\) test functions \(f_1,\ldots,f_m\) and
any fixed \((x,y)\), the matrix
\[
 M(x,y)=\bigl[\Gamma(f_i\otimes f_j)(x,y)\bigr]_{i,j=1}^m
\]
has a factorization \(U(x)\beta(x,y)V(y)^{\mathsf T}\), and hence
\begin{equation}\label{rank}
 \operatorname{rank} M(x,y)\leq \min(K_1,K_2).
\end{equation}

Choose \(m>\max(K_1,K_2)\). Take a nonempty open arc \(O\) of the
circle whose complement contains an open arc. In that complementary
arc choose \(m\) disjoint smaller arcs, and smooth real bump functions
\(f_i\), one supported in each smaller arc, normalized by
\(\int_\T f_i^2\,d\pi=1\).
Then \(f_i=0\) on \(O\) and \(\int_\T f_if_j\,d\pi=\delta_{ij}\).
Equation \eqref{bilinear} therefore yields
\[
 M(x,y)=I_m\qquad\text{for every }(x,y)\in O\times O,
\]
contradicting \eqref{rank}.
The contradiction also works if the factors are interpreted as
\(L^2\) representatives and the identities hold almost everywhere:
only finitely many factors and test functions occur, so choose a point
of the positive-measure set \(O\times O\) outside their exceptional
null sets. This completes the disproof.
\end{proof}

\newpage
\paragraph{Scope and verification.}
The final source permits locally compact Polish state spaces with a
smooth structure; the compact circle satisfies this condition.
Conjecture 19 is stated for ergodic Feller processes generally.
The following Proposition 20 treats the narrower continuous-path
diffusion case, which this jump-process example does not contradict.
The conclusion concerns the finite index ranges displayed in
equations (34), (37), and (39); an infinite expansion or an
auxiliary noise-space integral would be a different construction.
The accompanying exact checks verify projection, generator, marginal,
and rank identities in finite reset models. The arbitrary-rank
argument on smooth functions above supplies the full proof.

\paragraph{Task 2.}
Formulate and prove an exhaustion theorem that permits countably many
operators or integration over an auxiliary noise space, with explicit
convergence and domain conditions.

\paragraph{Task 3.}
Characterize jump kernels whose coupling operators admit a finite
product representation of the form \eqref{construction}.

\paragraph{Task 4.}
Determine the minimum numbers of marginal factors needed to represent
all coupling operators for finite-state irreducible Markov chains.

\paragraph{Task 5.}
Develop finite-rank approximations of shared-reset couplings and
quantify their effect on invariant measures and asymptotic variances.

\begin{thebibliography}{9}
\raggedright
\bibitem{NP}
N. N\"usken and G. A. Pavliotis,
\emph{Constructing Sampling Schemes via Coupling: Markov Semigroups and
Optimal Transport}, SIAM/ASA J. Uncertainty Quantification
\textbf{7}(1) (2019), 324--382,
\href{https://doi.org/10.1137/18M119896X}{doi:10.1137/18M119896X}.
Conjecture 19, equations (34)--(39), pp.~337--338;
Definition 6, p.~332; general setting, pp.~330--331.
\end{thebibliography}
\end{document}
